\section{投资问题：国内链能否接住文莱利润回落}

国内PTA、聚酯和CPL--PA6决定恒逸在文莱价差正常化后能否维持利润。核心判断是，PTA在集中检修和去库后出现加工费修复，聚酯销量具有韧性但库存与织机开工仍弱，CPL--PA6处于产能爬坡而非成熟盈利阶段。国内链提供次要改善，却尚不足以把文莱峰值利润替换为结构性增长。

2025年公司PTA收入44.16亿元、毛利率仅0.37\%；涤纶收入597.44亿元、毛利率4.07\%；锦纶及副产品收入18.20亿元、毛利率6.83\%。这组低毛利基线意味着小幅加工费变化可以带来较大利润弹性，也意味着行业产量增长并不自动转化为股东回报。

\section{PTA：检修修复有证据，长期过剩未消失}

行业口径显示，2025年中国大陆PTA产能约9,209万吨／年、产量约7,376万吨、表观需求约6,996万吨，出口约382万吨；平均加工差255.60元／吨，同比下降27.3\%，处于近五年低位。2026年6月，行业代理显示PTA开工约66.39\%、加工费约665元／吨、工厂利润约265元／吨，总库存约320.37万吨，集中检修推动阶段性修复。

但加工费上升的来源是供给检修而非终端强劲。若装置重启、聚酯需求没有跟上，库存和加工费可能再次逆转。恒逸2,150万吨参控股PTA产能分布在控股与参股平台，2025年报披露的控股口径产量和销量只有220.91／222.94万吨；因此不能用2,150万吨乘665元／吨推导公司利润。

\begin{exhibitbox}[PTA修复的证据与边界]
\noindent\begin{tabularx}{\textwidth}{L{3.1cm}R{2.2cm}L{2.5cm}X}
\toprule
\textbf{指标} & \textbf{行业／公司值} & \textbf{方向} & \textbf{投资边界} \\
\midrule
2025行业产能 & 9,209万吨／年 & 供给基数大 & 行业值，不是恒逸并表产能 \\
2025平均加工差 & 255.60元／吨 & 近五年低位 & 解释公司PTA低毛利方向 \\
2026-06加工费／利润 & 665／约265元／吨 & 检修后修复 & 阶段代理，不宜年化 \\
2026-06行业开工 & 66.39\% & 同比明显下降 & 不等于恒逸装置开工 \\
恒逸参控股产能 & 2,150万吨／年 & 规模领先 & 权益与合并口径混合 \\
恒逸2025控股产／销量 & 220.91／222.94万吨 & 可核验经营量 & 只覆盖浙江逸盛等控股口径 \\
恒逸2025PTA毛利率 & 0.37\% & 接近盈亏边缘 & 2026转正需半年报确认 \\
\bottomrule
\end{tabularx}
\end{exhibitbox}
\sourcenote{HY-OFF-2025AR；IND-PTA-202606；TONGKUN-2025引用CCF；行业与公司资料截止2026-07-22。行业加工费、开工和库存不等于恒逸单吨利润。}

投资上，PTA升级需要“加工费高于2025均值、行业不再累库、公司PTA毛利转正”同时出现；若加工费回落到约250元／吨以下或复产造成再累库，国内PTA不能接住文莱利润回落。

\section{聚酯：量增有韧性，库存决定价格传导}

恒逸2025年涤纶产量867.97万吨、销量868.93万吨，销量同比增长7.22\%，期末库存50.76万吨。收入／销量粗算均价约6,876元／吨，但该合并口径包含不同纤维和切片，不能视为统一ASP。聚酯纤维产能938万吨／年与产销量范围也未完全对齐，报告因此不发布分品种单吨利润。

行业需求呈现“产量增长、终端偏弱”。2025年化纤产量增长，聚酯产量约增长4\%，直纺长丝平均负荷约92\%；但加弹和织造平均开工约74\%／63\%，均同比下降。2026年6月行业代理显示聚酯开工约80.19\%、织机开工约57\%，POY／DTY／FDY库存约26.5／38.3／30.7天。上游检修带来原料修复，下游尚未形成全面补库。

\begin{exhibitbox}[聚酯验证仪表盘]
\small
\noindent\begin{tabularx}{\textwidth}{L{2.9cm}R{2.0cm}L{2.35cm}L{2.4cm}X}
\toprule
\textbf{指标} & \textbf{最新锚} & \textbf{当前判断} & \textbf{升级阈值} & \textbf{降级阈值} \\
\midrule
公司2025涤纶毛利率 & 4.07\% & 毛利薄 & 半年报毛利明显改善 & 销量增而毛利不升 \\
公司2025销量／库存 & 868.93／50.76万吨 & 量有韧性 & 库存增速低于销量 & 库存与短债同步升 \\
行业聚酯开工 & 80.19\% & 同比偏弱 & 开工回升且加工费扩大 & 继续下滑 \\
行业织机开工 & 57\% & 终端承接弱 & 回升并伴随库存下降 & 下滑且出口转弱 \\
POY／DTY／FDY库存 & 26.5／38.3／30.7天 & 库存偏高 & 三类库存持续下降 & 继续累库 \\
终端零售 & 2025约+3.2\% & 低个位数增长 & 零售／出口共同改善 & 贸易壁垒或内需转弱 \\
\bottomrule
\end{tabularx}
\end{exhibitbox}
\sourcenote{HY-OFF-2025AR；IND-PTA-202606；CCFA-2025；CNTAC-2025；行业锚截止2026-06，研究截止2026-07-22。}

聚酯主要采取库存销售和日度报价，不能套用制造业式backlog逻辑。未披露长约／现货占比、具名客户和分品种认证之前，最有效的领先指标是加工费、开工、库存和现金回款。只有量价库存同步改善，涤纶才能获得高于合并周期估值的信用。

\section{瓶片与RPET：产能存在，独立盈利证据不足}

公司披露聚酯瓶片含RPET参控股产能530万吨／年，需求锚为饮料、食品和包装。行业可用的理论价差为“瓶片价格减0.855倍PTA和0.335倍MEG”，但公司没有单列瓶片收入、实际开工、食品级认证、客户或ASP。产能数字能够证明业务直接性，却不能证明产销和利润。

瓶片新增供给、出口竞争、海运和食品接触认证都可能改变价格与市场准入。没有公司级证据时，本报告不给瓶片单独估值信用，也不把包装终端企业写成恒逸客户。若后续披露分部收入、认证和稳定毛利，可从“链条覆盖”升级为“可估值分部”。

\section{CPL--PA6：质量合格不等于盈利成熟}

广西项目使公司参控股CPL／PA6产能达到100／60万吨／年。年报称一期已全流程打通并进入试生产，主要产品质量合格率保持100\%；H1预告称项目贡献改善。与此同时，2025年锦纶及副产品产量8.89万吨、销量4.67万吨、库存4.22万吨，销量与设计产能之间仍有明显爬坡距离。

质量合格率是装置和产品可生产的证据，却不能替代良率、稳定负荷、客户认证、分产品ASP和单吨利润。工程塑料与高端锦纶通常需要更长认证和配方验证周期；在客户与认证未披露时，需求应用只作方向锚。

\begin{exhibitbox}[CPL--PA6从产能到盈利的升级路径]
\noindent\begin{tabularx}{\textwidth}{L{2.5cm}L{2.4cm}L{2.65cm}X}
\toprule
\textbf{阶段} & \textbf{现有证据} & \textbf{缺失字段} & \textbf{估值处理} \\
\midrule
装置建成／试生产 & 一期全流程打通 & 稳态负荷与检修 & 确认资产存在 \\
产品可用 & 公司称合格率100\% & 良率、等级与认证 & 不等于商业化份额 \\
销量爬坡 & 2025销4.67万吨 & CPL／PA6拆分与客户 & 观察收入路径 \\
毛利形成 & 锦纶及副产品毛利率6.83\% & 分产品单吨利润 & 仅合并分部信用 \\
稳定现金 & H1预告称改善 & 回款、库存与资本开支 & 基准不给成长溢价 \\
\bottomrule
\end{tabularx}
\end{exhibitbox}
\sourcenote{HY-OFF-2025AR；HY-OFF-2026H1P；analysis/value\_chain\_economics.md，资料截止2026-07-22。}

投资含义是，CPL--PA6可以作为事件期权，但不进入77.5亿元基准的额外增量瀑布。若半年报披露分产品产销量、良率、客户认证和正现金毛利，才允许提高2027--2028盈利或牛情景权重。

\section{终端需求：出口是边际，内需不是高增长}

2025年服装鞋帽针纺零售超过1.5万亿元、同比约增长3.2\%，网络服装零售约增长1.9\%。化纤出口约762万吨、同比增长14.64\%，长丝出口约435万吨、同比增长10.62\%。出口是重要边际需求，但也暴露于贸易壁垒、汇率和海运；内需低个位数增长不足以支撑聚酯长期双位数需求假设。

需求锚的正确用途是解释行业开工和库存，而不是证明恒逸订单。公司前五客户合计263.41亿元、占23.20\%，均未具名；重大销售合同不适用。只要客户和产品映射缺失，本报告不对客户集中度、长协占比或backlog给出溢价。

\begin{keyinsight}[国内链的“所以呢”]
国内链正在从低谷修复，但证据更像“亏损收窄／毛利回升”，不是“全面景气”。\textbf{行动上，PTA看加工费与库存，聚酯看三类长丝库存和织机开工，CPL--PA6看分产品产销与认证；三者没有同步改善前，不上调House 77.5亿元基准。}
\end{keyinsight}

\begin{riskbox}[国内链失效条件]
若PTA加工费重回约250元／吨以下、POY／DTY／FDY继续累库、织机开工低于57\%，且CPL--PA6库存接近或超过销量，则国内链无法接棒文莱利润。\textbf{行动上，下调2027--2028盈利恢复路径，并维持或扩大资本与现金流折价。}
\end{riskbox}
